\begin{thebibliography}{99}
\addcontentsline{toc}{section}{References}
\bibitem{wilson}
K.~G. Wilson, \emph{Confinement of quarks},
Physical Review D \textbf{10} (1974), 2445--2459.
\href{https://doi.org/10.1103/PhysRevD.10.2445}{doi:10.1103/PhysRevD.10.2445}.

\bibitem{luscher}
M. L\"uscher, \emph{Construction of a selfadjoint, strictly
positive transfer matrix for Euclidean lattice gauge theories},
Communications in Mathematical Physics \textbf{54} (1977), 283--292.
\href{https://doi.org/10.1007/BF01614090}{doi:10.1007/BF01614090}.
The finite-slab gluing argument used here is given directly
in Section~\ref{sec:boundary}.

\bibitem{cc}
A. Connes and C. Consani, \emph{Weil positivity and Trace
formula, the archimedean place}, 2020.
\href{https://arxiv.org/abs/2006.13771}{arXiv:2006.13771},
especially Appendices B--C for the explicit formula and
restricted positivity criterion. No short-support
positivity result is promoted here to a global one.

\bibitem{burnol}
J.-F. Burnol, \emph{Spectral Analysis of the local Conductor
Operator}, 1998.
\href{https://arxiv.org/abs/math/9811040}{arXiv:math/9811040}.

\bibitem{dlmf}
NIST Digital Library of Mathematical Functions,
\href{https://dlmf.nist.gov/5.5}{Section 5.5},
\href{https://dlmf.nist.gov/5.9}{Section 5.9},
\href{https://dlmf.nist.gov/5.11}{Section 5.11}
(gamma identities, digamma integrals and asymptotics);
\href{https://dlmf.nist.gov/25.10}{Section 25.10}
(zeta zeros);
\href{https://dlmf.nist.gov/27.12}{Section 27.12}
(prime number theorem).

\bibitem{elkies}
N.~D. Elkies, \emph{The product formula for $\xi(s)$ and
$\zeta(s)$; vertical distribution of zeros}, Math 229 lecture notes.
\href{https://people.math.harvard.edu/~elkies/M229.15/zeta2.pdf}{Author's lecture notes}.
Only the zero-counting estimate is used here; completion
conventions are fixed independently in this manuscript.

\bibitem{jaffe-witten}
A. Jaffe and E. Witten, \emph{Quantum Yang--Mills theory},
in \emph{The Millennium Prize Problems}, Clay Mathematics
Institute/American Mathematical Society.
\href{https://www.claymath.org/library/monographs/MPPc.pdf}{Official problem volume}.
Continuum existence and a mass gap are allowed assumptions,
not results claimed in this work.

\bibitem{buchholz-wichmann}
D. Buchholz and E.~H. Wichmann, \emph{Causal independence
and the energy-level density of states in local quantum
field theory}, Communications in Mathematical Physics
\textbf{106} (1986), 321--344.
\href{https://doi.org/10.1007/BF01454978}{doi:10.1007/BF01454978}.

\bibitem{mori}
H. Mori, \emph{Transport, Collective Motion, and Brownian
Motion}, Progress of Theoretical Physics \textbf{33}
(1965), 423--455.
\href{https://doi.org/10.1143/PTP.33.423}{doi:10.1143/PTP.33.423}.

\bibitem{meyer}
R. Meyer, \emph{A spectral interpretation for the zeros of
the Riemann zeta function}.
\href{https://arxiv.org/abs/math/0412277}{arXiv:math/0412277},
version 3 (2013 revision of the 2004 preprint),
especially Theorem 4.1, Corollary 4.2, and Theorem 5.8.

\bibitem{bw}
J.~J. Bisognano and E.~H. Wichmann,
\emph{On the duality condition for a Hermitian scalar field},
Journal of Mathematical Physics \textbf{16} (1975), 985--1007.
\href{https://doi.org/10.1063/1.522605}{doi:10.1063/1.522605}.
Its applicability to a proposed YM net is a separate
hypothesis.

\bibitem{runkel-szegedy}
I. Runkel and L. Szegedy, \emph{Area-Dependent Quantum
Field Theory}, Communications in Mathematical Physics
\textbf{381} (2021), 83--117.
\href{https://doi.org/10.1007/s00220-020-03902-1}{doi:10.1007/s00220-020-03902-1},
Section 5, especially Proposition 5.7.

\bibitem{levitin-strohmaier}
M. Levitin and A. Strohmaier,
\emph{Computations of eigenvalues and resonances on
perturbed hyperbolic surfaces with cusps}, 2018.
\href{https://arxiv.org/abs/1812.05554}{arXiv:1812.05554},
Section 7.3.2.
\end{thebibliography}
